% Generated by tools/edn_latex.py from reports/edn.md.
% Do not edit: change reports/edn.md and run `make edn`.
\ednentry{
  id = {EDN-05},
  title = {Supported makes: minimum listing support per make},
  date = {2026-09-22},
  milestone = {M1: Project Inception},
  topic = {Problem Specification},
  participants = {@lukas2510},
  decision = {Only makes with at least 300 listings in the cleaned used-car data (counted after the \texttt{ES} holdout, before the split) are supported; the pipeline computes the list and the API rejects other makes. With the current data this gives 11 makes covering 98.6\,\% of the used listings after the \texttt{ES} holdout (EDN-03).},
  rationale = {\emph{Alternatives considered:}
\begin{itemize}[leftmargin=*, topsep=2pt]
\item \textbf{Option A (chosen): minimum support threshold, reject the rest.}
\newline Pros: honest and testable scope; no silent low-quality predictions; easy to state in the model card.
\newline Cons: smaller product (e.g. Ford, Renault, Opel, Kia are not supported).
\item \textbf{Option B: accept all 25 makes and flag rare ones as low confidence.}
\newline Pros: broader API.
\newline Cons: weak predictions for makes with a few dozen listings; the warning logic needs its own tests and thresholds anyway.
\item \textbf{Option C: accept all makes without any rule.}
\newline Pros: simplest.
\newline Cons: no clear scope statement; hidden weaknesses for rare makes.
\end{itemize}
\par
\emph{Why:} The dataset is skewed toward premium brands (EDN-01 accepts this as a scope limit). A threshold turns that limit into an explicit, testable rule. 300 keeps 11 makes and 98.5\,\% of the data (98.6\,\% after the \texttt{ES} holdout); 200 would add three makes (0.7\,\% of the data), 500 would drop Aston Martin and Volkswagen.},
  ai = {Information seeking, Alternative generation, Alternative assessment, Recommendation},
  response = {Accepted},
  assessment = {AI counted the listings per make after scoping and deduplication, compared thresholds of 200, 300 and 500, and recommended option A. Lukas accepted it because it states the dataset's brand skew openly instead of hiding it behind a warning.},
  aievidence = {Claude Code session on 2026-09-22 while working on issue \#2: prompt ``How do we handle the 25 makes, many of which have very few listings?'' with options A to C; Lukas chose option A.},
  evidence = {\href{https://github.com/mlops-2627q1-mds-upc/MLOPS_Recommenditos/blob/main/docs/docs/problem-spec.md}{Problem specification}, section 2; \href{https://github.com/mlops-2627q1-mds-upc/MLOPS_Recommenditos/issues/2}{issue \#2}, \href{https://github.com/mlops-2627q1-mds-upc/MLOPS_Recommenditos/pull/17}{PR \#17}.},
}
